\subsection{Central polynomial coefficients and the remaining nonflat case}

The scalar theorem leaves a precise larger alternative:
a perfect diagonal coefficient can have polynomial
differentials. Diagonal terms alone do not supply
centrality for ordinary Soergel convolution. The
following results construct genuine central operations
and further restrict the remaining quadratic problem.

Initially let
\[
 A=k[x,y],\qquad \deg x=\deg y=2,\qquad
 s_1(x,y)=(-x,y),\quad s_2(x,y)=(x,-y),
\]
over a characteristic-zero field, and put
\[
 S=A^{\langle s_1,s_2\rangle}=k[x^2,y^2],\qquad
 B_i=A\otimes_{A^{s_i}}A(1),\quad b_i=[B_i].
\]
Retain \(\delta=v+v^{-1}\).
Write \(A\langle e\rangle\) for a free module whose
displayed generator has internal degree \(e\), without
changing the ordinary Soergel shift convention.

A \emph{perfect diagonal coefficient} \(C\) is a
bounded complex of finite graded free \(A\)-modules
with equal left and right actions. Its homogeneous
differential may have polynomial entries.
The target remains
\begin{equation}\label{eq:centralCoefficientTarget}
 [X]=P=b_1b_2+(b_1-\delta)(b_2-\delta),\qquad
 \chi_v(C)=\delta^2,
\end{equation}
with an actual graded derived bimodule isomorphism
\begin{equation}\label{eq:centralCoefficientQuadratic}
 X\otimes_A^{\mathbf L}X
       \simeq C\otimes_A^{\mathbf L}X.
\end{equation}
A central structure would additionally supply natural
isomorphisms moving \(C\) past the ordinary Soergel
objects and their maps, with unit and tensor coherence.
The necessary flatness conditions below do not assume
this extra structure, and hence apply when it is supplied.

\begin{finding}[\fid{F59}, polynomial centrality and flatness gates]
\label{f:centralPolynomialCoefficientGate}
For a nonzero homogeneous polynomial \(f\) of positive
degree, the diagonal Koszul complex
\[
 K(f)=[A\langle\deg f\rangle\xrightarrow{f}A]
\]
in degrees \(-1,0\) commutes with the actual generator
\(B_{s_i}\) in the derived bimodule category if and
only if \(s_i(f)=f\). Simultaneous commutation with
both generators is therefore equivalent to \(f\in S\).

Every bounded finite-free complex descended from \(S\)
has an explicit coherent central structure on the
ordinary Soergel category and its bounded homotopy
category, including its actual morphisms.

In every standard ordinary type-A ambient ring covered
by Finding~\ref{f:allRankScalarClosure}, the target
\eqref{eq:centralCoefficientTarget}--\eqref{eq:centralCoefficientQuadratic}
is impossible if all cohomology of \(C\) is flat.
It is also impossible if all cohomology of \(X\) is
flat on the left: that condition forces every
cohomology module of \(C\) to be graded free, even
though no flatness of \(C\) was assumed.
\end{finding}

\begin{proof}
\par\smallskip\noindent\textbf{The exact one-polynomial test.}\quad
Take \(s=s_1\) and suppress the common harmless
internal shift of \(B_s\). Its bimodule is the ring
\begin{equation}\label{eq:centralCoefficientReflectionRing}
 R_s=A\otimes_{A^s}A
     =k[x,x',y]/(x^2-x'^2).
\end{equation}
The two tensor complexes have differentials
\(f(x,y)\) and \(f(x',y)\), respectively.
The ring \(R_s\) is free over each copy of \(A\),
so both multiplications are injective. Ordinary
tensor computes the derived tensor, and the only
cohomology is in degree zero:
\begin{equation}\label{eq:centralCoefficientAnnihilators}
 R_s/(f(x,y)),\qquad R_s/(f(x',y)).
\end{equation}

The left and right \(A\)-actions generate \(R_s\).
An isomorphism of these bimodules is thus an
\(R_s\)-module isomorphism and forces equality
of their principal annihilator ideals.
In particular \(f(x,y)=a f(x',y)\) for some
\(a\in R_s\). Both polynomials have the same
degree, and multiplication by \(f(x',y)\) is
injective. The homogeneous decomposition of \(a\)
therefore puts it in degree zero, namely in \(k\).
On the diagonal quotient \(x'=x\), the nonzero
polynomial \(f\) forces \(a=1\).
On the other graph quotient \(x'=-x\), the
resulting equality gives \(f(x,y)=f(-x,y)\).
Thus \(s(f)=f\).

Conversely, an \(s\)-invariant \(f\) belongs
to \(A^s\), so its two actions on
\eqref{eq:centralCoefficientReflectionRing} are
identical. The two tensor complexes agree under
the standard identifications. The same proof for
\(s_2\) gives the assertion, and simultaneous
invariance is precisely \(f\in k[x^2,y^2]\).
This test is for the single Koszul complex; it is
not a classification of arbitrary differential
matrices with additional descent data.

For example \(f=x\) gives the two modules
\[
 k[x',y]/(x'^2),\qquad k[x,y]/(x^2).
\]
Left \(x\) is zero on the first and nonzero
nilpotent on the second. They are not isomorphic
even without grading.
On the graph module \(A_s\), however, the two
differentials are \(x\) and \(-x\), and multiplying
the degree-minus-one term by \(-1\) gives an
isomorphism. Over the fraction field both complexes
are acyclic. Graph commutation and generic checks
therefore do not supply actual generator centrality.

\par\smallskip\noindent\textbf{Coherent invariant descent.}\quad
Let \(K\) be a bounded complex of finite graded
free \(S\)-modules and put
\[
                         C_K=A\otimes_S K
\]
with its diagonal action. In this rank-one product,
\(A\) is free over \(S\) on \(1,x,y,xy\).
Every element of \(S\) has equal left and right
actions on \(B_i\), since \(S\subseteq A^{s_i}\).
This equality survives tensor products, sums,
shifts, direct summands and bimodule maps, hence
holds compatibly with every differential in an
ordinary Soergel complex \(Y\).

There are actual natural chain identifications
\begin{equation}\label{eq:centralCoefficientMovement}
 C_K\otimes_A Y\simeq K\otimes_S Y,
 \qquad
 Y\otimes_A C_K\simeq Y\otimes_S K.
\end{equation}
The first sends \((a\otimes k)\otimes y\) to
\(k\otimes ay\); the second sends
\(y\otimes(a\otimes k)\) to \(ya\otimes k\).
The balanced relations use the equal invariant
actions. The symmetry of complexes over the
commutative ring \(S\) sends
\(k\otimes y\) to \((-1)^{pq}y\otimes k\), where
\(p,q\) are their cohomological degrees.
It is \(A\)-bimodule-linear under these identifications,
commutes with differentials and is natural in all
ordinary Soergel maps.

Moving a labeled coefficient factor through an
iterated product gives the same permutation and
the same sum of Koszul exponents on either route.
The unit and tensor half-braiding identities are
therefore the ordinary symmetry/associativity
identities for complexes over \(S\).
This supplies actual coherent central maps even
when cohomology is not flat.
The same argument works with \(S=A^W\) in any
standard ordinary type-A product realization:
full invariants act equally on every generator.

\par\smallskip\noindent\textbf{Flat coefficient cohomology.}\quad
A finite graded flat module over the polynomial
ring \(A\) is finitely presented and projective;
this is the known equivalence in Stacks
Lemma~10.78.2, tag~00NX, parts~(1)--(3)
\cite{StacksFiniteProjective}.
It is graded free. Indeed lift a homogeneous
origin basis to a finite free surjection, split
by projectivity and use graded Nakayama on the
complementary kernel. The degree-zero component
of an ordinary splitting is a graded splitting.

If every \(H^j(C)\) is flat, successive
projective-cohomology splitting gives
\begin{equation}\label{eq:centralCoefficientScalarReduction}
 C\simeq\bigoplus_jH^j(C)[-j]
      \simeq C_{\mathrm{sc}}\otimes_k A,
\end{equation}
where \(C_{\mathrm{sc}}\) is a bounded finite
graded scalar complex with zero differential.
All these maps are \(A\)-linear between diagonal
objects, hence are also bimodule-linear.
Their Euler class is \(\chi_v(C_{\mathrm{sc}})=\delta^2\).
Substitution in
\eqref{eq:centralCoefficientQuadratic} gives the
scalar quadratic excluded by
Finding~\ref{f:allRankScalarClosure}.
A different central structure on \(C\) does not
alter this derived isomorphism.

\par\smallskip\noindent\textbf{Flat generator cohomology, without flatness of \(C\).}\quad
Let \(Q\) be the ungraded fraction field.
Since \(C\) is diagonal, \(C_Q\) is a scalar
complex of finite \(Q\)-vector spaces.
Let \(d(u)\) be its actual cohomological
dimension polynomial and \(F(u)\) the actual
generic graph character of \(X\).
The Euler classes and actual quadratic give
\[
 F^2=dF,\qquad F(-1)=2(e+\tau),\qquad d(-1)=4,
                  \quad\tau=(s,s).
\]
The positive idempotent-measure argument in the
preceding all-rank proof therefore gives
\begin{equation}\label{eq:centralCoefficientGenericGraphs}
                      F(u)=\frac{d(u)}2(e+\tau).
\end{equation}
This statement retains no internal grading.

Assume every \(M_i=H^i(X)\) is flat on the left.
It is then graded free. Its generic graph support
is contained in \(e,\tau\), so the codimension-two
graph lemma splits it globally into shifted
\(A_e,A_\tau\). In the actual ambient coordinates,
\(\tau\) negates two independent root variables
and fixes the others. The proof uses
\(yU=xV\), \(U^2=x^2I\), coprimality and the
two involution idempotents, with fixed coordinates
acting by their left scalars.
Every \(M_i\) is consequently right-free too.

The canonical cross-product K\"unneth map gives
\begin{equation}\label{eq:centralCoefficientFreeSquare}
 H^t(X\otimes_A^{\mathbf L}X)
       \simeq\bigoplus_{i+h=t}M_i\otimes_A M_h.
\end{equation}
It is justified by the already proved one-sided
projective-cohomology splitting, not by bimodule
formality. Its terms are free graph modules, so
this is a graded free left module in every degree.

For the right side of the quadratic, forget only
the \emph{right} action of the second factor \(X\).
Split its projective left cohomology in the
left-module complex category. Each \(M_i\) is
left-flat, so tensoring with the right complex
\(C\) gives, without any flatness assumption
on \(H^j(C)\), the actual left-module isomorphism
\begin{equation}\label{eq:centralCoefficientNonflatKunneth}
 H^t(C\otimes_A^{\mathbf L}X)
       \simeq\bigoplus_{j+i=t}H^j(C)\otimes_A M_i.
\end{equation}
The surviving left action is on \(C\); it
commutes with the chosen left-linear maps in
the second factor. Equivalently this is the
canonical cycle cross-product, whose bijectivity
is proved by the auxiliary one-sided splitting.
No bimodule splitting of \(X\) is asserted.

Choose a nonzero \(M_{i_0}\), which exists by
the nonzero generic Euler class. For diagonal
\(N=H^j(C)\), tensoring \(N\) on the right
with a graph module changes only its right
action and internal shift. Thus, in degree
\(t=j+i_0\), the direct sum in
\eqref{eq:centralCoefficientNonflatKunneth}
contains a shifted copy of \(H^j(C)\) as a
direct summand.
The actual quadratic identifies it with the
graded free left module
\eqref{eq:centralCoefficientFreeSquare}.
Therefore every \(H^j(C)\) is projective,
hence graded free. The scalar reduction above
gives the contradiction.
This proves the implication for nonflat \(C\)
as well as for already flat coefficients.
\end{proof}

In particular the earlier coupled Rouquier object
\[
 X_{\mathrm R}\simeq A_e^{\oplus2}
              \oplus A_\tau(-2)\oplus A_\tau(2)
\]
fails for \emph{every} perfect diagonal coefficient
of the specified Euler class.
Any remaining target construction must have some
nonflat cohomology on both sides, and must separately
provide its central structure if that operation is
required.

\subsubsection{An actual central quadratic with a different generator class}

Non-scalar central coefficients do occur in bounded
ordinary quadratics. The following explicit example
retains the required coefficient Euler class and
shows why the generator class remains essential.

In \(S\), put \(q=x^2\), of internal degree four,
and define
\[
 K_S(q)=[S\langle4\rangle\xrightarrow{q}S],\qquad
 U_S=S\oplus K_S(q)\oplus K_S(q)[1].
\]
The two Koszul Euler classes cancel, so
\(\chi_v(U_S)=1\). Since \(q\) is regular,
\[
 H^0(U_S)=S\oplus S/(q),\qquad
 H^{-1}(U_S)=S/(q),
\]
and all other groups vanish.
Let \(U=A\otimes_S U_S\). Its nonzero torsion
cohomology is \(A/(x^2)\), so it is not a scalar
complex of the kind excluded earlier.
Invariant descent supplies its coherent central maps.

Let \(B=B_1\otimes_A B_2\), and let
\(D_{\mathrm{sc}}\) be the cohomological-degree-zero
graded scalar space of character \(\delta^2\).
Define
\begin{equation}\label{eq:centralCoefficientDifferentClass}
 X'=U\otimes_A B,\qquad C'=D_{\mathrm{sc}}\otimes_k U.
\end{equation}
These are bounded ordinary Soergel and diagonal
coefficient complexes, respectively.
The rank-one ordinary decompositions give
\[
 B\otimes_A B\simeq D_{\mathrm{sc}}\otimes_k B.
\]
Indeed \(A\) is free over each reflection invariant
ring on \(1\) and the corresponding root, giving
the two shifted copies encoded by \(\delta\).
The two independent reflection factors commute,
giving \(\delta^2\); this is the ordinary
normalization in \cite[\S\S3.3--3.4, pp.~1105--1107]{EW}.

The actual central maps for \(U\) now give
\begin{align}\label{eq:centralCoefficientActualExample}
 X'\otimes_A X'
 &\simeq U\otimes_A U\otimes_A(B\otimes_A B)\notag\\
 &\simeq D_{\mathrm{sc}}\otimes_k U\otimes_A U\otimes_A B\notag\\
 &=C'\otimes_A X'.
\end{align}
These are actual chain/homotopy tensor isomorphisms.
The bounded terms are free on the relevant sides,
so they give the derived isomorphism too.
The central comparisons have the usual labeled
coefficient and Koszul coherence; no strict
identity of differently ordered raw reductions
is being claimed.

The Euler classes are
\[
                     \chi_v(C')=\delta^2,
                     \qquad [X']=b_1b_2.
\]
The latter is not the target \(P\): after forgetting
internal shifts it gives \(e+s_1+s_2+\tau\), whereas
\(P\) gives \(2e+2\tau\).
Both \(C'\) and \(X'\) have nonflat cohomology.
Thus the example proves a real central non-scalar
operation exists, while preserving the essential
ambient-class hypothesis of the exclusion.

\paragraph{The exact remaining gap.}
The nonflat/nonflat case for the designated \(P\)
is not decided here. One cannot specialize invariant
parameters at the origin and reuse the scalar theorem:
the resulting coinvariant quotient is not the
regular polynomial ring in its Tor argument.
The coefficient can also acquire torsion cohomology
or larger origin fibers.

For a scalar coefficient, its highest cohomology
survives generically; the global and generic top
comparisons therefore force the top module of \(X\)
to have positive rank and the required scalar
multiplicity. With a nonflat perfect diagonal
coefficient, an extra torsion top may vanish
generically. Those comparisons no longer identify
the two highest degrees. The results above force
nonflatness on both sides but do not supply this
missing identification.

The singular correspondence
\(A\otimes_{A^\tau}A\) is not thereby promoted to
a derived groupoid or a scalar coefficient.
Ordinary tensor reassociations do not prove a
derived quadratic when nonflat intersections
retain Tor. No such unverified comparison is
used in the construction or the exclusions.

% Separate paragraph for Part III.
\paragraph{The F59 central polynomial coefficients in Grothendieck's language.}
The diagonal module, reflection correspondence and
principal annihilator ideals are actual affine data.
EGA~I, Theorem~1.3.7, p.~86, Corollary~1.3.8,
p.~87 and Theorem~1.4.1, pp.~90--92 give their
coherent sheaf dictionary \cite{EGA}.
The Koszul test distinguishes two pullbacks to a
ramified correspondence; agreement on generic
graph pieces does not identify those pullbacks.
Invariant descent supplies an explicit pullback
from the invariant quotient with actual coherent
commuting operations, proved by balanced tensors.

The perfect diagonal complexes and derived tensors
fit SGA~6, I, Definition~4.7, pp.~121--122 and
Example~4.8, p.~122 \cite{SGA6}.
The flatness gate retains the actual canonical
K\"unneth maps and projective direct summands;
it does not use only their Euler classes.
The known flat/projective equivalence was cited
separately and its graded consequence was proved.

The grading and actual/virtual coefficient language
is SGA~3, I, 4.7.2--3, with cote~90, scan~17,
and SGA~6, V, \S\S2--3, with cote~133, scan~18
\cite{SGA3,SGA6,Archive}, retaining the previously
stated source-access qualification.
No \(\lambda\)-ring is assigned to the noncommutative
word algebra. The ordinary source normalization
and graph localization are
\cite[\S\S3.3--3.5]{EW} and
\cite[\S\S1.1--1.2, Theorem~1.6 and Remark~1.7]{EWLocalized}.
Those framework sources are distinct from the new
polynomial-coefficient constructions and deductions.
